\documentclass{beamer}
\hypersetup{pdfpagemode=FullScreen}
\setbeamertemplate{navigation symbols}{}
\usepackage[english]{babel}
\usepackage[latin1]{inputenc}
\usepackage{amsmath,amssymb,amsthm}
\usepackage[T1]{fontenc}
\usepackage{times}
\usepackage{mathptmx}
\usepackage{graphicx}
\usepackage{booktabs}
\usepackage{hyperref}
\usepackage{xcolor}
\usepackage{mathtools}
\usepackage{algorithm,algorithmic}
\usepackage{bm}

\renewcommand{\algorithmicrequire}{\textbf{Input:}}
\renewcommand{\algorithmicensure}{\textbf{Output:}}

\mode<presentation>{
  \usetheme{Madrid}
  \usecolortheme{whale}
}

\setbeamertemplate{blocks}[rounded][shadow=false]
\AtBeginEnvironment{block}{\small}
\AtBeginEnvironment{alertblock}{\small}
\AtBeginEnvironment{exampleblock}{\small}
\setbeamersize{text margin left=6pt, text margin right=6pt}
\setbeamercovered{transparent}
\setbeamertemplate{bibliography item}[text]
\setbeamertemplate{theorems}[numbered]
\setbeamerfont{title}{size=\large}
\setbeamerfont{date}{size=\tiny}
\setbeamertemplate{itemize items}[ball]
\setbeamertemplate{itemize subitem}[circle]
\setbeamertemplate{itemize subsubitem}[triangle]
\setbeamertemplate{caption}[numbered]

%------------------------------------------------------------------
%  TITLE PAGE
%------------------------------------------------------------------
\title[SVD \& PCA from Scratch]{%
  Singular Value Decomposition\\[4pt]
  and PCA from Scratch}

\author[Vatsav Reddy]{%
  \texorpdfstring{%
    {\Large Vatsav Reddy}\\[6pt]
    {\small Under the Guidance of}\\[2pt]
    {\small Prof.\ Praveen Parchuri}%
  }{Vatsav Reddy}}

\institute[IIITH]{%
  {\normalsize International Institute of Information Technology, Hyderabad}}

\date{\today}

\begin{document}

\begin{frame}
  \titlepage
\end{frame}

\begin{frame}[shrink=10]{Contents}
  \tableofcontents
\end{frame}

\AtBeginSection[]{
  \begin{frame}<beamer>
    \frametitle{Current Section}
    \tableofcontents[currentsection]
  \end{frame}
}

%==================================================================
\section{Variance, Covariance, and the Covariance Matrix}
%==================================================================

\begin{frame}[shrink=10]{Variance: Measuring Spread of a Single Variable}
  \begin{block}{Population Variance}
    For a random variable $x$ with mean $\mu_x = \mathbb{E}[x]$, the
    \textbf{variance} measures how far values typically deviate from the mean:
    \[
      \mathrm{Var}(x)
        = \mathbb{E}\!\left[(x - \mu_x)^2\right]
        = \mathbb{E}[x^2] - \left(\mathbb{E}[x]\right)^2
    \]
    The second equality follows by expanding $(x-\mu_x)^2$ and using
    linearity of expectation:
    $\mathbb{E}[(x-\mu_x)^2]
    = \mathbb{E}[x^2] - 2\mu_x\mathbb{E}[x] + \mu_x^2
    = \mathbb{E}[x^2] - \mu_x^2$.
  \end{block}

  \medskip

  \begin{block}{Sample Variance (from data)}
    Given $n$ observations $x_1, x_2, \ldots, x_n$ with sample mean
    $\bar{x} = \frac{1}{n}\sum_{i=1}^n x_i$:
    \[
      s^2 = \frac{1}{n-1}\sum_{i=1}^{n}(x_i - \bar{x})^2
    \]
    The factor $n-1$ (Bessel's correction) makes $s^2$ an unbiased
    estimator of the true population variance.
    In machine learning we often use $\frac{1}{n}$ for convenience.
  \end{block}

  \begin{alertblock}{Intuition}
    Large variance $\Rightarrow$ data is spread out; small variance
    $\Rightarrow$ data is tightly clustered near the mean.
    PCA will later exploit variance to identify \emph{informative} directions.
  \end{alertblock}
\end{frame}

%------------------------------------------------------------------
\begin{frame}[shrink=10]{Covariance: How Two Variables Move Together}
  \begin{block}{Population Covariance}
    For two random variables $x$ and $y$ with means $\mu_x, \mu_y$:
    \[
      \mathrm{Cov}(x,y)
        = \mathbb{E}\!\left[(x-\mu_x)(y-\mu_y)\right]
        = \mathbb{E}[xy] - \mathbb{E}[x]\,\mathbb{E}[y]
    \]
    This measures \emph{joint} variability.
  \end{block}

  \medskip

  \begin{block}{Sample Covariance}
    \[
      \widehat{\mathrm{Cov}}(x,y)
        = \frac{1}{n-1}\sum_{i=1}^{n}(x_i - \bar{x})(y_i - \bar{y})
    \]
    \begin{itemize}\setlength\itemsep{4pt}
      \item \textbf{Positive} $\Rightarrow$ $x$ and $y$ tend to increase
            \emph{together} (positive association)
      \item \textbf{Negative} $\Rightarrow$ one increases while the other
            decreases (opposite movement)
      \item \textbf{Near zero} $\Rightarrow$ no linear relationship between
            $x$ and $y$
    \end{itemize}
    Note: $\mathrm{Cov}(x,x) = \mathrm{Var}(x)$ --- covariance with
    itself is variance.
  \end{block}
\end{frame}

%------------------------------------------------------------------
\begin{frame}[shrink=10]{Correlation: Normalised Covariance}
  \begin{block}{Pearson Correlation Coefficient}
    Raw covariance depends on the \emph{units} of $x$ and $y$.
    Normalise by the product of standard deviations to get a
    unit-free measure:
    \[
      \rho_{xy}
        = \frac{\mathrm{Cov}(x,y)}{\sigma_x\,\sigma_y},
      \qquad
      \sigma_x = \sqrt{\mathrm{Var}(x)},\quad
      \sigma_y = \sqrt{\mathrm{Var}(y)}
    \]
    \[
      -1 \;\leq\; \rho_{xy} \;\leq\; 1
    \]
    \textbf{Sample version:}
    \[
      r_{xy}
        = \frac{\displaystyle\sum_{i=1}^n(x_i-\bar{x})(y_i-\bar{y})}
               {\displaystyle\sqrt{\sum_{i=1}^n(x_i-\bar{x})^2}
                \,\sqrt{\sum_{i=1}^n(y_i-\bar{y})^2}}
    \]
  \end{block}

  \begin{exampleblock}{Interpretation}
    $\rho = +1$: perfect positive linear relationship.
    $\rho = 0$: no linear relationship.
    $\rho = -1$: perfect negative linear relationship.
    High $|\rho|$ implies one variable is approximately a linear
    function of the other --- one is \emph{redundant}.
    PCA eliminates this redundancy.
  \end{exampleblock}
\end{frame}

%------------------------------------------------------------------
\begin{frame}[shrink=10]{The Covariance Matrix}
  \begin{block}{Setup: Multivariate Data}
    Let $\mathbf{x} = [x_1, x_2, \ldots, x_n]^\top \in \mathbb{R}^n$
    be a random vector.
    The \textbf{covariance matrix} $\Sigma \in \mathbb{R}^{n \times n}$ is:
    \[
      \Sigma = \mathbb{E}\!\left[
        (\mathbf{x}-\boldsymbol{\mu})(\mathbf{x}-\boldsymbol{\mu})^\top
      \right],
      \qquad
      \boldsymbol{\mu} = \mathbb{E}[\mathbf{x}]
    \]
    Entry $(i,j)$ of $\Sigma$ is $\Sigma_{ij} = \mathrm{Cov}(x_i, x_j)$.
  \end{block}

  \medskip

  \begin{block}{Explicit $2\times2$ Example}
    For $\mathbf{x} = [x, y]^\top$:
    \[
      \Sigma
        = \begin{bmatrix}
            \mathrm{Var}(x)   & \mathrm{Cov}(x,y)\\
            \mathrm{Cov}(y,x) & \mathrm{Var}(y)
          \end{bmatrix}
    \]
    The matrix is \textbf{symmetric} ($\Sigma_{ij} = \Sigma_{ji}$) because
    $\mathrm{Cov}(x_i,x_j) = \mathrm{Cov}(x_j,x_i)$.
  \end{block}

  \medskip

  \begin{block}{Correlation Matrix}
    Normalise each entry by $\sigma_i\sigma_j$:
    $R_{ij} = \dfrac{\mathrm{Cov}(x_i,x_j)}{\sigma_i\sigma_j}$,\quad
    so $R_{ii} = 1$ (diagonal is all ones).
  \end{block}
\end{frame}

%------------------------------------------------------------------
\begin{frame}[shrink=10]{Sample Covariance Matrix from Data}
  \begin{block}{Data Matrix and Sample Covariance}
    With $m$ data points in $\mathbb{R}^n$, arrange them as rows of
    $X \in \mathbb{R}^{m \times n}$.
    After centring (zero-mean columns):
    $X_c = X - \mathbf{1}\boldsymbol{\mu}^\top$
    where $\boldsymbol{\mu} = \frac{1}{m}\sum_{i=1}^m \mathbf{x}_i$.

    \medskip

    The \textbf{sample covariance matrix} is:
    \[
      \boxed{
        \Sigma
          = \frac{1}{m}\,X_c^\top X_c
          \;\in\; \mathbb{R}^{n \times n}
      }
    \]
    Entry $(j,k)$:
    $\Sigma_{jk}
       = \frac{1}{m}\sum_{i=1}^{m}(X_c)_{ij}(X_c)_{ik}$
    $=$ sample covariance between feature $j$ and feature $k$.
  \end{block}

  \begin{exampleblock}{Two Key Properties}
    \begin{enumerate}\setlength\itemsep{5pt}
      \item \textbf{Symmetric}: $\Sigma^\top = \frac{1}{m}(X_c^\top X_c)^\top
            = \frac{1}{m}X_c^\top X_c = \Sigma$.\quad
            Follows from $(AB)^\top = B^\top A^\top$.
      \item \textbf{Positive semi-definite (PSD)}:
            For any $\mathbf{v} \in \mathbb{R}^n$,\;
            $\mathbf{v}^\top \Sigma\, \mathbf{v}
             = \frac{1}{m}\mathbf{v}^\top X_c^\top X_c\mathbf{v}
             = \frac{1}{m}\|X_c\mathbf{v}\|^2 \geq 0$.
            Hence all eigenvalues $\lambda_i \geq 0$.
    \end{enumerate}
  \end{exampleblock}
\end{frame}

%==================================================================
\section{Eigenvalue Decomposition (EVD) and Spectral Theorem}
%==================================================================

\begin{frame}[shrink=10]{Eigenvectors and Eigenvalues: Recap}
  \begin{block}{Definition}
    For a square matrix $A \in \mathbb{R}^{n \times n}$, a nonzero
    vector $\mathbf{v}$ is an \textbf{eigenvector} with
    \textbf{eigenvalue} $\lambda$ if:
    \[
      A\mathbf{v} = \lambda\mathbf{v}
    \]
    $A$ stretches or flips $\mathbf{v}$ along its own direction ---
    \emph{no rotation}, only scaling.
    To find $\lambda$: rearrange as $(A - \lambda I)\mathbf{v} = \mathbf{0}$.
    For a nonzero solution $\mathbf{v}$ to exist, the matrix
    $(A-\lambda I)$ must be singular, i.e.\ $\det(A-\lambda I) = 0$.
    This is the \emph{characteristic polynomial} of $A$.
  \end{block}

  \medskip

  \begin{block}{Why Eigenvectors Matter: Matrix Multiplication Becomes Scalar}
    If $\mathbf{x} = \sum_{i=1}^n \alpha_i \mathbf{v}_i$ (eigenvectors
    as basis), then:
    \[
      A\mathbf{x}
        = A\sum_{i=1}^n \alpha_i\mathbf{v}_i
        = \sum_{i=1}^n \alpha_i A\mathbf{v}_i
        = \sum_{i=1}^n \alpha_i \lambda_i \mathbf{v}_i
    \]
    The matrix multiplication reduces to \textbf{scalar multiplications}
    $\alpha_i \mapsto \alpha_i\lambda_i$.
    This is the fundamental computational advantage of the eigenbasis.
  \end{block}
\end{frame}

%------------------------------------------------------------------
\begin{frame}[shrink=10]{Eigenvalue Decomposition (EVD)}
  \begin{block}{Construction}
    Suppose $A$ has $n$ linearly independent eigenvectors
    $\mathbf{v}_1, \ldots, \mathbf{v}_n$ with eigenvalues
    $\lambda_1, \ldots, \lambda_n$.
    Stack eigenvectors as columns of $V \in \mathbb{R}^{n \times n}$:
    \[
      AV
        = A[\mathbf{v}_1 \;\cdots\; \mathbf{v}_n]
        = [\lambda_1\mathbf{v}_1 \;\cdots\; \lambda_n\mathbf{v}_n]
        = V\Lambda
    \]
    where $\Lambda = \mathrm{diag}(\lambda_1,\ldots,\lambda_n)$.
    If $V$ is invertible (eigenvectors linearly independent):
    \[
      \boxed{A = V\Lambda V^{-1}}
      \qquad\Longleftrightarrow\qquad
      \Lambda = V^{-1}AV
    \]
  \end{block}

  \medskip

  \begin{block}{Spectral Theorem for Symmetric Matrices}
    If $A$ is \textbf{real symmetric} ($A = A^\top$), then:
    \begin{enumerate}\setlength\itemsep{4pt}
      \item All eigenvalues are \textbf{real}.
      \item Eigenvectors corresponding to \emph{distinct} eigenvalues
            are \textbf{orthogonal}.
      \item $A$ has a complete orthonormal set of eigenvectors, so
            $V$ can be chosen to satisfy $V^\top V = I$,
            meaning $V^{-1} = V^\top$.
    \end{enumerate}
    Therefore, for a real symmetric $A$:
    \[
      \boxed{A = V\Lambda V^\top}
    \]
    This is the \textbf{Spectral Decomposition}.
  \end{block}
\end{frame}

%------------------------------------------------------------------
\begin{frame}[shrink=10]{Applying EVD to the Covariance Matrix}
  \begin{block}{Diagonalising $\Sigma = \frac{1}{m}X_c^\top X_c$}
    Since $\Sigma$ is real symmetric and PSD, the Spectral Theorem gives:
    \[
      \Sigma\,\mathbf{p}_i = \lambda_i\,\mathbf{p}_i,
      \quad
      \lambda_1 \geq \lambda_2 \geq \cdots \geq \lambda_n \geq 0,
      \quad
      \mathbf{p}_i^\top\mathbf{p}_j = \delta_{ij}
    \]
    Collecting into matrices:
    \[
      \Sigma = P\Lambda P^\top,
      \quad
      P = [\mathbf{p}_1 \;\cdots\; \mathbf{p}_n],
      \quad
      \Lambda = \mathrm{diag}(\lambda_1,\ldots,\lambda_n)
    \]
    \textbf{Key insight:} $\mathbf{p}_i^\top \Sigma\, \mathbf{p}_i = \lambda_i$,
    so the variance of the projected data along direction $\mathbf{p}_i$
    equals $\lambda_i$. Large eigenvalue $=$ high variance direction.
  \end{block}

  \begin{alertblock}{Computational Note}
    In practice, use \texttt{numpy.linalg.eigh} (not \texttt{eig}) for
    symmetric matrices: it returns \emph{real} eigenvalues in ascending
    order and is more numerically stable.
    Reverse to get descending order (largest first).
  \end{alertblock}
\end{frame}

%==================================================================
\section{The Gram Matrix Trick: Avoiding $O(n^3)$ Cost}
%==================================================================

\begin{frame}[shrink=10]{The Computational Problem}
  \begin{block}{Why Direct EVD of $\Sigma$ is Expensive}
    For face images of size $100 \times 100$, we have $n = 10{,}000$
    features (pixels).
    The covariance matrix $\Sigma = \frac{1}{m}X_c^\top X_c$ is then
    $n \times n = 10{,}000 \times 10{,}000$.
    \begin{itemize}\setlength\itemsep{5pt}
      \item \textbf{Memory:} $n^2 = 10^8$ floating point entries
            $\approx 800\,\text{MB}$
      \item \textbf{Eigendecomposition:} $O(n^3) = O(10^{12})$
            operations --- impractical
    \end{itemize}
  \end{block}

  \medskip

  \begin{block}{The Gram Matrix Alternative}
    Instead of diagonalising the $n \times n$ matrix $\Sigma$,
    diagonalise the much smaller $m \times m$
    \textbf{Gram matrix}:
    \[
      G = \frac{1}{m}\,X_c X_c^\top \;\in\; \mathbb{R}^{m \times m}
    \]
    With $m = 400$ images and $n = 10{,}000$ pixels, this gives an EVD
    cost of $O(m^3) = O(400^3) = O(6.4 \times 10^7)$ operations instead
    of $O(10^{12})$ --- a speedup of
    $\left(\frac{n}{m}\right)^3 = 25^3 = 15{,}625\times$.
  \end{block}
\end{frame}

%------------------------------------------------------------------
\begin{frame}[shrink=10]{Why $G$ and $\Sigma$ Share the Same Non-Zero Eigenvalues}
  \begin{block}{Shared Spectrum Identity}
    Let $G\mathbf{v}_i = \lambda_i\mathbf{v}_i$ for some nonzero
    $\mathbf{v}_i \in \mathbb{R}^m$ and $\lambda_i \neq 0$.
    We claim $\mathbf{p}_i := X_c^\top\mathbf{v}_i$ is an eigenvector
    of $\Sigma$ with the \emph{same} eigenvalue $\lambda_i$.

    \medskip

    \textbf{Proof:} Expand $G\mathbf{v}_i = \lambda_i\mathbf{v}_i$:
    \[
      \frac{1}{m}\,X_c X_c^\top\mathbf{v}_i = \lambda_i\mathbf{v}_i
    \]
    Multiply both sides on the left by $X_c^\top$:
    \[
      X_c^\top\!\left(\frac{1}{m}\,X_c X_c^\top\mathbf{v}_i\right)
        = X_c^\top(\lambda_i\mathbf{v}_i)
    \]
    \[
      \underbrace{\frac{1}{m}\,X_c^\top X_c}_{\Sigma}
        \underbrace{X_c^\top\mathbf{v}_i}_{\mathbf{p}_i}
        = \lambda_i \underbrace{X_c^\top\mathbf{v}_i}_{\mathbf{p}_i}
    \]
    \[
      \therefore\quad
      \boxed{\Sigma\,\mathbf{p}_i = \lambda_i\,\mathbf{p}_i}
    \]
    $\mathbf{p}_i = X_c^\top\mathbf{v}_i$ is indeed an eigenvector of
    $\Sigma$ with the \emph{same} eigenvalue $\lambda_i$.
    \hfill $\square$
  \end{block}
\end{frame}

%------------------------------------------------------------------
\begin{frame}[shrink=10]{Normalising the Pixel-Space Eigenvector}
  \begin{block}{Deriving the Normalisation Factor $\sqrt{m\lambda_i}$}
    We have $\mathbf{p}_i \propto X_c^\top\mathbf{v}_i$.
    To make $\|\mathbf{p}_i\| = 1$, compute the norm:
    \[
      \|X_c^\top\mathbf{v}_i\|^2
        = (X_c^\top\mathbf{v}_i)^\top(X_c^\top\mathbf{v}_i)
        = \mathbf{v}_i^\top X_c X_c^\top \mathbf{v}_i
    \]
    Since $X_c X_c^\top = mG$ (by definition), and
    $G\mathbf{v}_i = \lambda_i\mathbf{v}_i$:
    \[
      = \mathbf{v}_i^\top (mG)\mathbf{v}_i
        = m\,\mathbf{v}_i^\top G\mathbf{v}_i
        = m\,\mathbf{v}_i^\top (\lambda_i\mathbf{v}_i)
        = m\lambda_i\,\underbrace{\|\mathbf{v}_i\|^2}_{=1}
        = m\lambda_i
    \]
    Therefore $\|X_c^\top\mathbf{v}_i\| = \sqrt{m\lambda_i}$, and
    the \textbf{normalised pixel-space eigenvector} (eigenface) is:
    \[
      \boxed{
        \mathbf{p}_i
          = \frac{X_c^\top\mathbf{v}_i}{\sqrt{m\lambda_i}}
      },
      \qquad
      \|\mathbf{p}_i\| = 1
    \]
  \end{block}

  \begin{exampleblock}{Interpretation}
    $\mathbf{v}_i \in \mathbb{R}^m$: coefficients telling how to
    linearly combine the $m$ training images.
    $X_c^\top\mathbf{v}_i$: the resulting pixel pattern from that
    combination.
    After normalisation: unit-norm direction in pixel space.
  \end{exampleblock}
\end{frame}

%==================================================================
\section{The Full PCA Pipeline}
%==================================================================

\begin{frame}[shrink=10]{PCA: Step-by-Step Algorithm}
  \begin{block}{Step 1 --- Centre the Data}
    Compute the mean image:
    $\boldsymbol{\mu} = \frac{1}{m}\sum_{i=1}^m \mathbf{x}_i
                      \in \mathbb{R}^n$.
    Centre:
    $X_c = X - \mathbf{1}_m\boldsymbol{\mu}^\top$.
    Now every column of $X_c$ has zero mean.
    \emph{Why?} Covariance measures spread \emph{around the mean};
    centring ensures $\Sigma = \frac{1}{m}X_c^\top X_c$ correctly
    represents variances and covariances.
  \end{block}

  \begin{block}{Step 2 --- Eigendecompose (via Gram Matrix if $m \ll n$)}
    Compute $G = \frac{1}{m}X_c X_c^\top$ and solve $G\mathbf{v}_i = \lambda_i\mathbf{v}_i$.
    Recover pixel-space eigenvectors via
    $\mathbf{p}_i = \frac{X_c^\top\mathbf{v}_i}{\sqrt{m\lambda_i}}$.
    Order by $\lambda_1 \geq \lambda_2 \geq \cdots$.
  \end{block}

  \begin{block}{Step 3 --- Select Top-$k$ Principal Components}
    Form $P_k = [\mathbf{p}_1,\ldots,\mathbf{p}_k]^\top \in \mathbb{R}^{k \times n}$
    (each row is a principal direction).
    These $k$ directions capture the most variance.
  \end{block}

  \begin{block}{Step 4 --- Project}
    \[
      Z = X_c P_k^\top \in \mathbb{R}^{m \times k}
    \]
    Each row $\mathbf{z}_i = X_{c,i}P_k^\top \in \mathbb{R}^k$ gives
    the PCA coefficients (scores) for data point $i$.
    The $j$-th score $z_{ij} = \mathbf{x}_{c,i}^\top \mathbf{p}_j$:
    the projection of image $i$ onto eigenface $j$.
  \end{block}
\end{frame}

%------------------------------------------------------------------
\begin{frame}[shrink=10]{Reconstruction and Reconstruction Error}
  \begin{block}{Reconstruction Formula}
    Recover the original (approximate) image from PCA codes:
    \[
      \hat{\mathbf{x}}_i
        = \boldsymbol{\mu} + \mathbf{z}_i P_k
        = \boldsymbol{\mu} + \mathbf{x}_{c,i} P_k^\top P_k
    \]
    In summation form (most interpretable):
    \[
      \boxed{
        \hat{\mathbf{x}}_i
          = \boldsymbol{\mu}
            + \sum_{j=1}^{k}
              \underbrace{\bigl(\mathbf{x}_{c,i}^\top\mathbf{p}_j\bigr)}_{z_{ij}}
              \mathbf{p}_j
      }
    \]
    \textbf{Reading this:}
    Start from the mean face $\boldsymbol{\mu}$.
    Compute how much of each eigenface $\mathbf{p}_j$ is present
    (the scalar $z_{ij}$).
    Add those weighted eigenfaces back.
  \end{block}

  \begin{block}{Why $P_k^\top P_k \neq I$}
    $P_k \in \mathbb{R}^{k \times n}$ has orthonormal \emph{rows},
    so $P_k P_k^\top = I_k$ (rows are unit, orthogonal).
    But $P_k^\top P_k \in \mathbb{R}^{n \times n}$ is a rank-$k$
    orthogonal projector onto the subspace spanned by
    $\mathbf{p}_1,\ldots,\mathbf{p}_k$, satisfying
    $(P_k^\top P_k)^2 = P_k^\top P_k$.
    The reconstruction projects $\mathbf{x}_{c,i}$ onto this
    $k$-dimensional subspace.
  \end{block}
\end{frame}

%------------------------------------------------------------------
\begin{frame}[shrink=10]{Explained Variance Ratio}
  \begin{block}{How Much Information Does Each Component Capture?}
    The variance of the projected data along direction $\mathbf{p}_i$:
    \[
      \frac{1}{m}\|\hat{X}_i\|^2
        = \frac{1}{m}(X_c\mathbf{p}_i)^\top(X_c\mathbf{p}_i)
        = \frac{1}{m}\mathbf{p}_i^\top X_c^\top X_c\mathbf{p}_i
        = \mathbf{p}_i^\top \Sigma\,\mathbf{p}_i
        = \mathbf{p}_i^\top (\lambda_i\mathbf{p}_i)
        = \lambda_i
    \]
    So the variance along the $i$-th principal component is exactly
    $\lambda_i$ (the corresponding eigenvalue).

    \medskip

    The \textbf{Explained Variance Ratio} (EVR) for component $i$:
    \[
      \mathrm{EVR}_i
        = \frac{\lambda_i}{\sum_{j=1}^n \lambda_j}
    \]
    This is the fraction of \emph{total variance} captured by
    direction $\mathbf{p}_i$.
    To choose $k$: find the smallest $k$ such that
    $\sum_{i=1}^k \mathrm{EVR}_i$ exceeds a threshold (e.g.\ $0.95$).
  \end{block}

  \begin{alertblock}{Practical Advice (from lecture)}
    Set aside 100 validation images.
    Compute reconstruction error for $k = 100, 1000, 10{,}000$.
    The right $k$ is application-dependent: medical/industrial images
    need much lower reconstruction error than coarse face recognition.
  \end{alertblock}
\end{frame}

%==================================================================
\section{Singular Value Decomposition (SVD)}
%==================================================================

\begin{frame}[shrink=10]{The Problem with Rectangular Matrices}
  \begin{block}{Eigenvectors Require Square Matrices}
    The eigenvalue equation $A\mathbf{v} = \lambda\mathbf{v}$ requires
    both sides to live in the \emph{same} space $\mathbb{R}^n$.
    For a \textbf{rectangular} matrix $A \in \mathbb{R}^{m \times n}$
    with $m \neq n$:
    \[
      A\underbrace{\mathbf{v}}_{\in\,\mathbb{R}^n}
        = \underbrace{\mathbf{w}}_{\in\,\mathbb{R}^m}
    \]
    The output lives in a \emph{different} space than the input.
    The equation $A\mathbf{v} = \lambda\mathbf{v}$ cannot hold because
    $\lambda\mathbf{v} \in \mathbb{R}^n$ while $A\mathbf{v} \in \mathbb{R}^m$.
    Hence rectangular matrices \textbf{cannot have eigenvectors}.
  \end{block}

  \begin{block}{Why This Is a Real Problem}
    Data matrices $X \in \mathbb{R}^{m \times n}$ are almost always
    rectangular.
    We \emph{want} the computational advantage of eigenvectors
    (converting matrix multiplication to scalar multiplication), but
    we cannot directly apply EVD.
    SVD is the resolution: it finds \emph{two} sets of orthonormal
    vectors --- one for the input space $\mathbb{R}^n$ and one for the
    output space $\mathbb{R}^m$ --- that \emph{together} make the
    matrix operation scalar.
  \end{block}
\end{frame}

%------------------------------------------------------------------
\begin{frame}[shrink=10]{SVD: The Core Idea}
  \begin{block}{Looking for Paired Bases in Two Different Spaces}
    Suppose we can find:
    \begin{itemize}\setlength\itemsep{4pt}
      \item $\mathbf{v}_1,\ldots,\mathbf{v}_k \in \mathbb{R}^n$
            (orthonormal, forming a basis for the \emph{row space} of $A$)
      \item $\mathbf{u}_1,\ldots,\mathbf{u}_k \in \mathbb{R}^m$
            (orthonormal, forming a basis for the \emph{column space} of $A$)
    \end{itemize}
    such that:
    \[
      A\mathbf{v}_i = \sigma_i\,\mathbf{u}_i,
      \quad i = 1,\ldots,k
    \]
    where $\sigma_i > 0$ are scalars called \textbf{singular values}.

    \medskip

    If every $\mathbf{x} \in \mathbb{R}^n$ is written as
    $\mathbf{x} = \sum_{i=1}^k \alpha_i\mathbf{v}_i$, then:
    \[
      A\mathbf{x}
        = \sum_{i=1}^k \alpha_i A\mathbf{v}_i
        = \sum_{i=1}^k \alpha_i\sigma_i\mathbf{u}_i
    \]
    Once again, \textbf{matrix multiplication reduces to scalar
    multiplication} --- this is the exact same advantage as eigenvectors,
    but now valid for rectangular matrices.
  \end{block}
\end{frame}

%------------------------------------------------------------------
\begin{frame}[shrink=10]{SVD: Matrix Form}
  \begin{block}{Writing $A\mathbf{v}_i = \sigma_i\mathbf{u}_i$ as a Matrix Equation}
    Stack all relations simultaneously.
    Let $V = [\mathbf{v}_1\;\cdots\;\mathbf{v}_n] \in \mathbb{R}^{n\times n}$
    (all right singular vectors, extended to a full basis of $\mathbb{R}^n$
    via Gram--Schmidt for the null space directions).
    Let $U = [\mathbf{u}_1\;\cdots\;\mathbf{u}_m] \in \mathbb{R}^{m\times m}$
    (left singular vectors, similarly extended).
    Let $\Sigma \in \mathbb{R}^{m\times n}$ be the matrix with
    $\Sigma_{ii} = \sigma_i$ for $i=1,\ldots,k$ and zeros elsewhere.
    Then:
    \[
      AV = U\Sigma
      \quad\Rightarrow\quad
      \boxed{A = U\Sigma V^\top}
    \]
    This is the \textbf{Singular Value Decomposition} of $A$.
    Since $U$ and $V$ are orthogonal ($UU^\top = I$, $VV^\top = I$),
    their inverses are their transposes.
  \end{block}

  \begin{block}{Dimensions}
    $U \in \mathbb{R}^{m\times m}$ (left singular vectors),\;
    $\Sigma \in \mathbb{R}^{m\times n}$ (diagonal, singular values),\;
    $V \in \mathbb{R}^{n\times n}$ (right singular vectors).
    The \emph{thin} (economy) SVD keeps only the $k = \mathrm{rank}(A)$
    non-zero singular values:
    $U_k \in \mathbb{R}^{m\times k}$,\;
    $\Sigma_k \in \mathbb{R}^{k\times k}$,\;
    $V_k \in \mathbb{R}^{n\times k}$.
  \end{block}
\end{frame}

%------------------------------------------------------------------
\begin{frame}[shrink=10]{How to Find $U$ and $V$: Connection to EVD}
  \begin{block}{Computing $V$ (Right Singular Vectors)}
    Suppose $A = U\Sigma V^\top$ exists.
    Compute $A^\top A$:
    \[
      A^\top A
        = (U\Sigma V^\top)^\top (U\Sigma V^\top)
        = V\Sigma^\top U^\top U\Sigma V^\top
        = V\underbrace{\Sigma^\top\Sigma}_{\text{diagonal: }\sigma_i^2}V^\top
    \]
    (using $U^\top U = I$ since $U$ is orthogonal).
    This is exactly the \textbf{EVD of $A^\top A$}:
    $A^\top A = V\Lambda_V V^\top$ where $\Lambda_V = \Sigma^\top\Sigma$
    has entries $\lambda_i = \sigma_i^2$.
    Therefore: \textbf{$V$ consists of the eigenvectors of $A^\top A$}.
  \end{block}

  \begin{block}{Computing $U$ (Left Singular Vectors)}
    Similarly, compute $AA^\top$:
    \[
      AA^\top
        = (U\Sigma V^\top)(U\Sigma V^\top)^\top
        = U\Sigma V^\top V\Sigma^\top U^\top
        = U\underbrace{\Sigma\Sigma^\top}_{\text{diagonal: }\sigma_i^2}U^\top
    \]
    This is the EVD of $AA^\top$:
    $AA^\top = U\Lambda_U U^\top$.
    Therefore: \textbf{$U$ consists of the eigenvectors of $AA^\top$}.
    And the singular values $\sigma_i = \sqrt{\lambda_i}$
    (square roots of the common eigenvalues of $A^\top A$ and $AA^\top$).
  \end{block}
\end{frame}

%------------------------------------------------------------------
\begin{frame}[shrink=10]{Why SVD Always Exists}
  \begin{block}{SVD vs.\ EVD: Existence Conditions}
    EVD of $A$ requires:
    \begin{enumerate}\setlength\itemsep{4pt}
      \item $A$ is \textbf{square} ($n \times n$)
      \item $A$ has $n$ linearly independent eigenvectors
            (not always guaranteed)
    \end{enumerate}
    SVD has \textbf{no such restrictions}. It always exists because:
    \begin{itemize}\setlength\itemsep{4pt}
      \item $A^\top A$ is always a \textbf{real symmetric PSD} matrix
            (regardless of the shape of $A$).
            Symmetric matrices always have a complete set of real
            orthogonal eigenvectors (Spectral Theorem).
      \item Similarly $AA^\top$ is always real symmetric PSD.
    \end{itemize}
    Therefore, $V$ and $U$ always exist, and hence \textbf{every real
    matrix $A \in \mathbb{R}^{m \times n}$ has an SVD}.
  \end{block}

  \begin{exampleblock}{Geometric Interpretation}
    Any linear transformation $\mathbb{R}^n \to \mathbb{R}^m$ can be
    decomposed into three steps:
    (1) rotate/reflect in the input space (by $V^\top$);
    (2) scale along axes (by $\Sigma$);
    (3) rotate/reflect in the output space (by $U$).
  \end{exampleblock}
\end{frame}

%==================================================================
\section{Low-Rank Approximation and the Eckart--Young Theorem}
%==================================================================

\begin{frame}[shrink=10]{SVD as a Sum of Rank-1 Matrices}
  \begin{block}{Outer Product Expansion}
    From $A = U\Sigma V^\top = \sum_{i=1}^k \sigma_i\mathbf{u}_i\mathbf{v}_i^\top$:
    \[
      A
        = \sigma_1\mathbf{u}_1\mathbf{v}_1^\top
        + \sigma_2\mathbf{u}_2\mathbf{v}_2^\top
        + \cdots
        + \sigma_k\mathbf{u}_k\mathbf{v}_k^\top
    \]
    Each term $\sigma_i\mathbf{u}_i\mathbf{v}_i^\top$ is a
    \textbf{rank-1 matrix} (outer product of two vectors).
    They are added in order of decreasing importance (since
    $\sigma_1 \geq \sigma_2 \geq \cdots \geq \sigma_k \geq 0$).
    Including all $k$ terms recovers $A$ exactly.
    \textbf{Dropping the last $k - r$ terms} gives a rank-$r$
    approximation.
  \end{block}

  \begin{block}{Storage Comparison}
    Original $A$: $m \times n$ entries.
    Rank-$r$ approximation: need to store $r$ triples
    $(\mathbf{u}_i, \sigma_i, \mathbf{v}_i)$:
    $(m + n + 1) \times r$ values.
    Compression ratio $\approx \dfrac{mn}{(m+n)r}$.
    For $m=400$, $n=10{,}000$, $r=50$:
    original $4 \times 10^6$, compressed $\approx 5 \times 10^5$
    --- 8$\times$ reduction.
  \end{block}
\end{frame}

%------------------------------------------------------------------
\begin{frame}[shrink=10]{The Eckart--Young--Mirsky Theorem}
  \begin{theorem}[Best Low-Rank Approximation]
    Let $A \in \mathbb{R}^{m \times n}$ with SVD
    $A = \sum_{i=1}^k \sigma_i\mathbf{u}_i\mathbf{v}_i^\top$.
    The \textbf{best rank-$r$ approximation} of $A$ in the Frobenius
    norm is:
    \[
      A_r = \sum_{i=1}^r \sigma_i\mathbf{u}_i\mathbf{v}_i^\top
    \]
    That is, $A_r$ solves:
    \[
      \min_{B\,:\,\mathrm{rank}(B) \leq r}
        \|A - B\|_F^2
    \]
    and the minimum error is:
    \[
      \|A - A_r\|_F^2 = \sum_{i=r+1}^k \sigma_i^2
    \]
  \end{theorem}

  \begin{block}{Why This Is Deep}
    The theorem says: if you must approximate $A$ by a rank-$r$ matrix
    and you want to minimise the total squared entry-wise error,
    the answer is exactly the first $r$ terms of the SVD expansion.
    \textbf{The singular values directly quantify reconstruction quality.}
  \end{block}
\end{frame}

%------------------------------------------------------------------
\begin{frame}[shrink=10]{Frobenius Norm: What We Are Minimising}
  \begin{block}{Definition}
    The Frobenius norm of a matrix $A \in \mathbb{R}^{m \times n}$:
    \[
      \|A\|_F = \sqrt{\sum_{i=1}^m\sum_{j=1}^n a_{ij}^2}
              = \sqrt{\mathrm{tr}(A^\top A)}
    \]
    It is the \emph{square root of the sum of squared entries} ---
    the natural generalisation of the Euclidean norm to matrices.
    $\|A - B\|_F^2 = \sum_{i,j}(a_{ij} - b_{ij})^2$ measures
    entry-wise squared differences.
  \end{block}

  \begin{block}{Connection to Singular Values}
    Since $\mathrm{tr}(A^\top A) = \mathrm{tr}(V\Sigma^2 V^\top)
    = \mathrm{tr}(\Sigma^2) = \sum_i \sigma_i^2$:
    \[
      \|A\|_F^2 = \sum_{i=1}^k \sigma_i^2
    \]
    So the singular values fully determine the Frobenius norm.
    Dropping the $r+1,\ldots,k$ terms loses
    $\sum_{i=r+1}^k \sigma_i^2$ --- the ``energy'' in those directions.
  \end{block}
\end{frame}

%==================================================================
\section{SVD and PCA: The Complete Connection}
%==================================================================

\begin{frame}[shrink=10]{PCA via SVD of the Centred Data Matrix}
  \begin{block}{Direct SVD of $X_c$}
    Apply SVD to the centred data matrix $X_c \in \mathbb{R}^{m\times n}$:
    \[
      X_c = U\Sigma V^\top
    \]
    $V \in \mathbb{R}^{n \times n}$: right singular vectors.

    Now compute $X_c^\top X_c$:
    \[
      X_c^\top X_c
        = (U\Sigma V^\top)^\top(U\Sigma V^\top)
        = V\Sigma^\top U^\top U\Sigma V^\top
        = V\Sigma^2 V^\top
    \]
    Comparing with the EVD $\Sigma_{\mathrm{cov}} = P\Lambda P^\top$:
    \begin{itemize}\setlength\itemsep{4pt}
      \item \textbf{Principal directions} $P = V$ (right singular vectors of $X_c$)
      \item \textbf{Eigenvalues} $\Lambda = \frac{1}{m}\Sigma^2$,\quad
            i.e.\ $\lambda_i = \sigma_i^2/m$
    \end{itemize}
    Therefore, \textbf{PCA of $X_c$ is equivalent to the SVD of $X_c$}.
    One can run PCA by directly SVD-decomposing the data matrix ---
    often more numerically stable than computing $X_c^\top X_c$ first.
  \end{block}
\end{frame}

%------------------------------------------------------------------
\begin{frame}[shrink=10]{Row Space, Column Space, Rank, and Null Space}
  \begin{block}{Fundamental Subspaces of $A \in \mathbb{R}^{m \times n}$}
    \begin{itemize}\setlength\itemsep{5pt}
      \item \textbf{Column space} (range of $A$):
            $\mathrm{col}(A) = \{A\mathbf{x} : \mathbf{x} \in \mathbb{R}^n\}
            \subseteq \mathbb{R}^m$.
            Dimension $= \mathrm{rank}(A)$.
      \item \textbf{Row space}: $\mathrm{col}(A^\top) \subseteq \mathbb{R}^n$.
            Dimension $= \mathrm{rank}(A)$.
      \item \textbf{Null space}: $\mathrm{null}(A) = \{\mathbf{x} : A\mathbf{x} = \mathbf{0}\}
            \subseteq \mathbb{R}^n$.
            Dimension $= n - \mathrm{rank}(A)$ (Rank--Nullity Theorem).
      \item \textbf{Left null space}: $\mathrm{null}(A^\top) \subseteq \mathbb{R}^m$.
    \end{itemize}
  \end{block}

  \begin{block}{SVD Reveals the Subspaces}
    The right singular vectors $\mathbf{v}_1,\ldots,\mathbf{v}_k$ span the
    \textbf{row space} of $A$.
    The remaining $\mathbf{v}_{k+1},\ldots,\mathbf{v}_n$ span the
    \textbf{null space}.
    The left singular vectors $\mathbf{u}_1,\ldots,\mathbf{u}_k$ span the
    \textbf{column space}.
    The remaining $\mathbf{u}_{k+1},\ldots,\mathbf{u}_m$ span the
    \textbf{left null space}.

    \medskip

    \textbf{Consequence}: only the $k$ dimensional subspace of $\mathbb{R}^n$
    in the row space can generate a non-zero output; null-space vectors
    are annihilated.
    So we only need $k = \mathrm{rank}(A)$ basis vectors to fully
    represent what $A$ does.
  \end{block}
\end{frame}

%==================================================================
\section{Eigenfaces: PCA Applied to Face Recognition}
%==================================================================

\begin{frame}[shrink=10]{Eigenfaces: Setup}
  \begin{block}{The Data}
    $m$ training face images, each of size $100 \times 100$ pixels
    $\Rightarrow$ $n = 10{,}000$ features per image.
    Data matrix $X \in \mathbb{R}^{m \times 10{,}000}$.
    After centring: $X_c$ with zero-mean columns.
  \end{block}

  \begin{block}{Computing Eigenfaces via the Gram Matrix Trick}
    Direct EVD of $\Sigma$ is infeasible ($10{,}000 \times 10{,}000$).
    Instead:
    \begin{enumerate}\setlength\itemsep{5pt}
      \item Compute $G = \frac{1}{m}X_c X_c^\top \in \mathbb{R}^{m\times m}$
      \item Solve $G\mathbf{v}_i = \lambda_i\mathbf{v}_i$ (cheap: $m \times m$)
      \item Recover eigenfaces:
            $\mathbf{p}_i = \frac{X_c^\top\mathbf{v}_i}{\sqrt{m\lambda_i}}$
            $\in \mathbb{R}^{10{,}000}$
      \item Reshape each $\mathbf{p}_i$ to $100 \times 100$ for
            visualisation --- these are the \textbf{eigenfaces}
      \item Keep top-$k$ eigenfaces (e.g.\ $k = 50$--$100$)
    \end{enumerate}
  \end{block}

  \begin{block}{Properties of Eigenfaces}
    \begin{itemize}\setlength\itemsep{4pt}
      \item They form an orthonormal basis for the space of faces in
            the dataset
      \item Any face $\approx \boldsymbol{\mu} + \sum_{i=1}^k z_i\mathbf{p}_i$
      \item Store $\mathbf{p}_1,\ldots,\mathbf{p}_k$ once;
            then only $k$ scalars $z_i$ per new image
    \end{itemize}
  \end{block}
\end{frame}

%------------------------------------------------------------------
\begin{frame}[shrink=10]{Eigenfaces: Compression and Recognition}
  \begin{block}{Storage Analysis}
    Original dataset: $m \times n = m \times 10{,}000$ numbers.
    After PCA with $k$ components:
    \[
      \underbrace{k \times n}_{\text{eigenfaces }P_k}
      + \underbrace{n}_{\text{mean face }\boldsymbol{\mu}}
      + \underbrace{m \times k}_{\text{scores }Z}
      \;\text{ numbers total}
    \]
    Numerical example: $m = 400$, $n = 10{,}000$, $k = 50$:
    \begin{itemize}\setlength\itemsep{3pt}
      \item Original: $400 \times 10{,}000 = 4{,}000{,}000$ values
      \item Compressed: $50 \times 10{,}000 + 10{,}000 + 400 \times 50
            = 510{,}000$ values
      \item Compression factor: $\approx 7.8\times$
    \end{itemize}
  \end{block}

  \begin{block}{Generalisation and Limitations}
    The eigenfaces basis is learned from training data.
    Any new face can be projected onto this basis and (approximately)
    reconstructed.
    If training data is diverse (many demographics, lighting conditions),
    the basis generalises well.
    If training data is narrow (only one demographic), the basis will
    reconstruct out-of-distribution faces poorly.
    Same principle as any machine learning model: quality of the learned
    representation depends on diversity of training data.
  \end{block}
\end{frame}

%==================================================================
\section{Summary: SVD, PCA, and Their Connections}
%==================================================================

\begin{frame}[shrink=10]{Summary: The Big Picture}
  \begin{columns}[T]
    \column{0.48\textwidth}
    \begin{block}{EVD (Square Matrices Only)}
      $A = V\Lambda V^{-1}$\\
      Symmetric: $A = V\Lambda V^\top$\\
      \textbf{Requires:} square; sufficient independent eigenvectors\\
      \textbf{Gives:} eigenvalues $\lambda_i$, eigenvectors $\mathbf{v}_i$\\
      $A\mathbf{v}_i = \lambda_i\mathbf{v}_i$
    \end{block}

    \begin{block}{SVD (Any Matrix)}
      $A = U\Sigma V^\top$\\
      $U$: left singular vectors (eigenvectors of $AA^\top$)\\
      $V$: right singular vectors (eigenvectors of $A^\top A$)\\
      $\sigma_i = \sqrt{\lambda_i}$: singular values\\
      Always exists; no shape restriction
    \end{block}

    \column{0.48\textwidth}
    \begin{block}{PCA (via EVD or SVD)}
      \textbf{Goal:} Find low-dimensional representation\\
      \textbf{Tool:} EVD of $\Sigma = \frac{1}{m}X_c^\top X_c$\\
      \textbf{Shortcut:} SVD of $X_c$ directly\\
      \textbf{Principal directions} $=$ eigenvectors of $\Sigma$\\
            $=$ right singular vectors of $X_c$\\
      \textbf{Variance} along $\mathbf{p}_i$ $=$ $\lambda_i$
    \end{block}

    \begin{alertblock}{Gram Matrix Trick}
      When $m \ll n$:\\
      EVD of $G = \frac{1}{m}X_c X_c^\top$ ($m{\times}m$)\\
      then $\mathbf{p}_i = \frac{X_c^\top\mathbf{v}_i}{\sqrt{m\lambda_i}}$\\
      Cost: $O(m^3)$ instead of $O(n^3)$
    \end{alertblock}
  \end{columns}

  \medskip

  \begin{center}
    \large\textit{%
      ``SVD always exists. EVD needs square matrices.
        PCA is EVD of the covariance matrix --- or equivalently,
        SVD of the data matrix.''}
  \end{center}
\end{frame}

%==================================================================
\section{Summary of Basics}
%==================================================================

%------------------------------------------------------------------
\begin{frame}[shrink=10]{Cauchy--Schwarz Inequality}
  \begin{block}{Statement}
    For any two vectors $\mathbf{u}, \mathbf{v}$ in an inner product
    space:
    \[
      \lvert\langle\mathbf{u},\mathbf{v}\rangle\rvert
        \;\leq\; \|\mathbf{u}\|\,\|\mathbf{v}\|
    \]
    In $\mathbb{R}^n$ (dot product form):
    \[
      \left(\sum_{i=1}^n u_i v_i\right)^2
        \leq
        \left(\sum_{i=1}^n u_i^2\right)
        \left(\sum_{i=1}^n v_i^2\right)
    \]
  \end{block}

  \begin{block}{Proof via Cosine Formula}
    By definition of the inner product in $\mathbb{R}^n$:
    \[
      \langle\mathbf{u},\mathbf{v}\rangle
        = \|\mathbf{u}\|\,\|\mathbf{v}\|\cos\theta
    \]
    Since $|\cos\theta| \leq 1$ for all $\theta \in [0,\pi]$:
    \[
      |\langle\mathbf{u},\mathbf{v}\rangle|
        = \|\mathbf{u}\|\,\|\mathbf{v}\|\,|\cos\theta|
        \leq \|\mathbf{u}\|\,\|\mathbf{v}\|
    \]
    \textbf{Equality} holds iff $|\cos\theta| = 1$, i.e.\
    $\theta = 0$ or $\pi$, i.e.\
    $\mathbf{u} = \lambda\mathbf{v}$ for some scalar $\lambda$
    (vectors are linearly dependent).
  \end{block}

  \begin{exampleblock}{Connection to PCA}
    The projection coefficient $\alpha_i = \mathbf{z}^\top\mathbf{p}_i$
    satisfies $|\alpha_i| \leq \|\mathbf{z}\|\|\mathbf{p}_i\| = \|\mathbf{z}\|$
    (since $\|\mathbf{p}_i\|=1$). This bounds how much variance a single
    principal component can capture.
    The correlation coefficient $\rho_{xy} = \frac{\mathrm{Cov}(x,y)}{\sigma_x\sigma_y}$
    satisfies $|\rho_{xy}| \leq 1$ precisely because of Cauchy--Schwarz applied
    to the vectors of centred observations.
  \end{exampleblock}
\end{frame}

%------------------------------------------------------------------
\begin{frame}[shrink=10]{Numerical Example: EVD of $AA^\top$}
  \begin{block}{Setup}
    Let $A = \begin{bmatrix} 3 & 0\\ 4 & 5 \end{bmatrix}$.
    Compute:
    \[
      AA^\top
        = \begin{bmatrix}3&0\\4&5\end{bmatrix}
          \begin{bmatrix}3&4\\0&5\end{bmatrix}
        = \begin{bmatrix}9&12\\12&41\end{bmatrix}
    \]
  \end{block}

  \begin{block}{Finding Eigenvalues: Characteristic Polynomial}
    \[
      \det(AA^\top - \lambda I)
        = \det\begin{bmatrix}9-\lambda & 12\\ 12 & 41-\lambda\end{bmatrix}
        = (9-\lambda)(41-\lambda) - 144 = 0
    \]
    Expanding: $\lambda^2 - 50\lambda + 369 - 144 = \lambda^2 - 50\lambda + 225 = 0$.
    Quadratic formula:
    $\lambda = \dfrac{50 \pm \sqrt{2500 - 900}}{2} = \dfrac{50 \pm 40}{2}$.
    \[
      \lambda_1 = 45,\qquad \lambda_2 = 5
    \]
    Singular values of $A$: $\sigma_1 = \sqrt{45} = 3\sqrt{5}$,\;
    $\sigma_2 = \sqrt{5}$.
  \end{block}

  \begin{block}{Finding Eigenvectors}
    For $\lambda_1 = 45$: $(AA^\top - 45I)\mathbf{v} = \mathbf{0}$
    gives $-36x + 12y = 0 \Rightarrow y = 3x$.
    Normalised: $\mathbf{u}_1 = \frac{1}{\sqrt{10}}[1,\,3]^\top$.\\[4pt]
    For $\lambda_2 = 5$: $4x + 12y = 0 \Rightarrow x = -3y$.
    Normalised: $\mathbf{u}_2 = \frac{1}{\sqrt{10}}[-3,\,1]^\top$.\\[4pt]
    Check orthogonality: $\mathbf{u}_1^\top\mathbf{u}_2
    = \frac{1}{10}(1\cdot(-3) + 3\cdot 1) = 0$.\;\checkmark\quad
    (Guaranteed since $AA^\top$ is symmetric.)
  \end{block}
\end{frame}

%------------------------------------------------------------------
\begin{frame}[shrink=10]{Condition Number and Why SVD Is Numerically Superior}
  \begin{block}{Definition: Condition Number}
    For a matrix $X$ with singular values $\sigma_1 \geq \sigma_2 \geq \cdots$,
    the \textbf{condition number} is:
    \[
      \kappa(X) = \frac{\sigma_{\max}}{\sigma_{\min}}
    \]
    It measures how much the matrix can \emph{amplify} relative errors
    in the input.
    Large $\kappa$ means the matrix is nearly singular and numerical
    errors are magnified.
  \end{block}

  \begin{block}{Why Forming $X^\top X$ Squares the Condition Number}
    Let $X = U\Sigma V^\top$ (SVD). Then:
    \[
      X^\top X = V\Sigma^\top U^\top U\Sigma V^\top = V\Sigma^2 V^\top
    \]
    The eigenvalues of $X^\top X$ are $\sigma_i^2$ (squares of singular
    values of $X$). Therefore:
    \[
      \kappa(X^\top X)
        = \frac{\sigma_{\max}^2}{\sigma_{\min}^2}
        = \left(\frac{\sigma_{\max}}{\sigma_{\min}}\right)^2
        = \kappa(X)^2
    \]
    \textbf{Example:} If $\kappa(X) = 10^3$, then
    $\kappa(X^\top X) = 10^6$.
    If $\sigma_{\min} = 10^{-8}$, then $\sigma_{\min}^2 = 10^{-16}$,
    which is at the threshold of 64-bit floating point machine precision
    ($\approx 10^{-16}$) --- the computer may treat it as zero,
    completely destroying that direction's information.
  \end{block}

  \begin{alertblock}{Practical Conclusion}
    Direct SVD algorithms (Golub--Kahan bidiagonalisation, QR iterations
    with Householder reflections) compute $U, \Sigma, V$ from $X$
    \emph{without ever forming} $X^\top X$.
    They work with the original singular values $\sigma_i$ rather than
    their squares, keeping numerical errors minimal.
    This is why \texttt{sklearn.decomposition.PCA} uses SVD internally.
  \end{alertblock}
\end{frame}

%------------------------------------------------------------------
\begin{frame}[shrink=10]{Theory vs.\ Practice: Finding $V$ in SVD}
  \begin{block}{What Theory Says}
    From $A = U\Sigma V^\top$, compute:
    \[
      A^\top A = V\Sigma^2 V^\top
    \]
    So $V$ consists of the eigenvectors of $A^\top A$, and
    $U$ consists of the eigenvectors of $AA^\top$.
    The singular values satisfy $\sigma_i = \sqrt{\lambda_i}$.
    This is the mathematical \emph{derivation} that connects SVD to EVD.
  \end{block}

  \begin{block}{What Practice Does}
    Modern SVD algorithms \textbf{do not form} $A^\top A$ or $AA^\top$.
    They work directly on $A$ using orthogonal transformations:
    \begin{enumerate}\setlength\itemsep{4pt}
      \item \textbf{Bidiagonalisation} (Householder reflections):
            reduce $A$ to a bidiagonal matrix $B$ such that
            $A = U_1 B V_1^\top$ with orthogonal $U_1, V_1$.
      \item \textbf{QR iteration} on $B$: compute SVD of $B$
            numerically without forming $B^\top B$.
    \end{enumerate}
    At no step are eigenvalues of $A^\top A$ computed.
    The singular values are found directly as $\sigma_i$, not via
    $\sqrt{\lambda_i(A^\top A)}$, so the squaring step and its
    numerical damage are avoided entirely.
  \end{block}

  \begin{exampleblock}{The Key Distinction}
    \textbf{Theory:} $V$ are eigenvectors of $A^\top A$ --- useful for
    understanding and deriving.\\
    \textbf{Practice:} SVD algorithms compute $V$ from $A$ directly,
    bypassing $A^\top A$ to preserve numerical precision.
  \end{exampleblock}
\end{frame}

%------------------------------------------------------------------
\begin{frame}[shrink=10]{Bessel's Correction: Why Divide by $n-1$?}
  \begin{block}{The Problem with Dividing by $n$}
    Given i.i.d.\ samples $X_1,\ldots,X_n$ with true variance $\sigma^2$,
    define the naive sample variance:
    $S_n^2 = \frac{1}{n}\sum_{i=1}^n (X_i - \bar{X})^2$.
    We will show $\mathbb{E}[S_n^2] = \frac{n-1}{n}\sigma^2 \neq \sigma^2$
    --- it is \textbf{biased downward}.
  \end{block}

  \begin{block}{Key Identity}
    Expand by adding and subtracting $\mu$:
    \begin{align*}
      \sum_{i=1}^n(X_i - \bar{X})^2
        &= \sum_{i=1}^n\bigl[(X_i-\mu) - (\bar{X}-\mu)\bigr]^2\\
        &= \sum_{i=1}^n(X_i-\mu)^2
           - 2(\bar{X}-\mu)\sum_{i=1}^n(X_i-\mu)
           + n(\bar{X}-\mu)^2\\
        &= \sum_{i=1}^n(X_i-\mu)^2 - n(\bar{X}-\mu)^2
    \end{align*}
    where we used $\sum_{i=1}^n(X_i - \mu) = n\bar{X} - n\mu = n(\bar{X}-\mu)$.
  \end{block}
\end{frame}

%------------------------------------------------------------------
\begin{frame}[shrink=10]{Bessel's Correction: Full Derivation}
  \begin{block}{Taking Expectations}
    Apply $\mathbb{E}[\,\cdot\,]$ (using linearity) to the identity
    $\sum(X_i-\bar{X})^2 = \sum(X_i-\mu)^2 - n(\bar{X}-\mu)^2$:
    \[
      \mathbb{E}\!\left[\sum_{i=1}^n(X_i-\bar{X})^2\right]
        = \mathbb{E}\!\left[\sum_{i=1}^n(X_i-\mu)^2\right]
          - n\,\mathbb{E}\!\left[(\bar{X}-\mu)^2\right]
    \]

    \textbf{First term:}
    $\mathbb{E}[(X_i-\mu)^2] = \mathrm{Var}(X_i) = \sigma^2$
    for each $i$, so
    $\mathbb{E}\!\left[\sum_{i=1}^n(X_i-\mu)^2\right] = n\sigma^2$.

    \medskip

    \textbf{Second term:}
    $\mathbb{E}[(\bar{X}-\mu)^2] = \mathrm{Var}(\bar{X})$.
    Since $X_i$ are i.i.d.:
    \[
      \mathrm{Var}(\bar{X})
        = \mathrm{Var}\!\left(\frac{1}{n}\sum_{i=1}^n X_i\right)
        = \frac{1}{n^2}\sum_{i=1}^n\mathrm{Var}(X_i)
        = \frac{1}{n^2}\cdot n\sigma^2
        = \frac{\sigma^2}{n}
    \]
    So $n\,\mathbb{E}[(\bar{X}-\mu)^2] = n\cdot\frac{\sigma^2}{n} = \sigma^2$.

    \medskip

    \textbf{Combining:}
    \[
      \mathbb{E}\!\left[\sum_{i=1}^n(X_i-\bar{X})^2\right]
        = n\sigma^2 - \sigma^2 = (n-1)\sigma^2
    \]
  \end{block}
\end{frame}

%------------------------------------------------------------------
\begin{frame}[shrink=10]{Bessel's Correction: Conclusion and Degrees of Freedom}
  \begin{block}{Biased vs.\ Unbiased Estimator}
    From the derivation: $\mathbb{E}\!\left[\sum_{i=1}^n(X_i-\bar{X})^2\right] = (n-1)\sigma^2$.

    \medskip

    Dividing by $n$ gives a \textbf{biased} estimator:
    \[
      \mathbb{E}[S_n^2]
        = \mathbb{E}\!\left[\frac{1}{n}\sum(X_i-\bar{X})^2\right]
        = \frac{1}{n}(n-1)\sigma^2
        = \frac{n-1}{n}\,\sigma^2 \;\neq\; \sigma^2
    \]
    Dividing by $n-1$ gives an \textbf{unbiased} estimator:
    \[
      \mathbb{E}[S^2]
        = \mathbb{E}\!\left[\frac{1}{n-1}\sum(X_i-\bar{X})^2\right]
        = \frac{1}{n-1}(n-1)\sigma^2
        = \sigma^2
      \quad\checkmark
    \]
  \end{block}

  \begin{block}{Why $n-1$? The Degrees of Freedom Argument}
    The deviations $d_i = X_i - \bar{X}$ satisfy the constraint
    $\sum_{i=1}^n d_i = 0$ (deviations always sum to zero by definition
    of the mean).
    This means once $n-1$ deviations are known, the $n$-th is completely
    determined.
    Only $n-1$ deviations are \emph{free to vary} --- there are
    $n-1$ \textbf{degrees of freedom}.
    Intuitively: we used one degree of freedom to \emph{estimate} $\mu$
    by $\bar{X}$ from the same data, leaving $n-1$ independent pieces
    of information for the variance estimate.
  \end{block}

  \begin{alertblock}{Connection to PCA / Covariance Matrix}
    In PCA with zero-mean data, we use $\Sigma = \frac{1}{m}X_c^\top X_c$.
    The factor $m$ is a convention common in machine learning
    (population variance, not sample variance).
    If using sample covariance (statistics convention), divide by $m-1$.
    For large $m$, the difference is negligible and does not affect which
    eigenvectors are found --- only the scale of eigenvalues.
  \end{alertblock}
\end{frame}

%------------------------------------------------------------------
\begin{frame}
  \Huge{\centerline{Thank You}}
  \vspace{0.6cm}
  \large{\centerline{Questions?}}
\end{frame}

%==================================================================
\section*{Summary of Basics}
%==================================================================

%------------------------------------------------------------------
% (A) CAUCHY--SCHWARZ
%------------------------------------------------------------------
\begin{frame}[shrink=10]{Summary of Basics --- Cauchy--Schwarz Inequality}
  \begin{block}{Statement (Inner-Product Form)}
    For any vectors $\mathbf{u}, \mathbf{v}$ in an inner-product space:
    \[
      |\langle \mathbf{u}, \mathbf{v}\rangle|
        \;\leq\;
        \|\mathbf{u}\|\,\|\mathbf{v}\|
    \]
    In $\mathbb{R}^n$ with the standard dot product:
    \[
      \left(\sum_{i=1}^n u_i v_i\right)^{\!2}
        \;\leq\;
        \left(\sum_{i=1}^n u_i^2\right)
        \left(\sum_{i=1}^n v_i^2\right)
    \]
  \end{block}

  \begin{block}{Proof via the Cosine Formula}
    Recall the projection of $\mathbf{u}$ onto $\mathbf{v}$:
    \[
      \langle \mathbf{u},\mathbf{v}\rangle
        = \|\mathbf{u}\|\,\|\mathbf{v}\|\cos\theta
    \]
    where $\theta$ is the angle between them.
    Since $|\cos\theta| \leq 1$ for all $\theta$:
    \[
      |\langle \mathbf{u},\mathbf{v}\rangle|
        = \|\mathbf{u}\|\,\|\mathbf{v}\|\,|\cos\theta|
        \leq \|\mathbf{u}\|\,\|\mathbf{v}\|
    \]
    Equality holds $\Leftrightarrow |\cos\theta|=1 \Leftrightarrow
    \theta \in \{0^\circ,180^\circ\} \Leftrightarrow
    \mathbf{u} = \lambda\mathbf{v}$ for some scalar $\lambda$
    (vectors are \emph{linearly dependent}).
  \end{block}

  \begin{exampleblock}{Numerical Check: $\mathbf{u}=(1,2),\,\mathbf{v}=(3,4)$}
    $\langle\mathbf{u},\mathbf{v}\rangle = 11$;\quad
    $\|\mathbf{u}\|\|\mathbf{v}\| = \sqrt{5}\cdot 5 = 5\sqrt{5} \approx 11.18$.
    Indeed $11 \leq 11.18$. \checkmark
  \end{exampleblock}

  \begin{alertblock}{Why It Appears in PCA}
    The correlation coefficient $\rho_{xy} = \frac{\mathrm{Cov}(x,y)}{\sigma_x\sigma_y}$
    satisfies $|\rho_{xy}| \leq 1$ precisely because of Cauchy--Schwarz:
    $|\mathrm{Cov}(x,y)| \leq \sqrt{\mathrm{Var}(x)\,\mathrm{Var}(y)}$.
    This guarantees the covariance matrix is PSD and that correlation
    is a bounded, well-defined measure.
  \end{alertblock}
\end{frame}

%------------------------------------------------------------------
% (B) CONDITION NUMBER
%------------------------------------------------------------------
\begin{frame}[shrink=10]{Summary of Basics --- Condition Number}
  \begin{block}{Definition}
    For a matrix $X$ with singular values
    $\sigma_1 \geq \sigma_2 \geq \cdots \geq \sigma_r > 0$, the
    \textbf{condition number} is:
    \[
      \kappa(X)
        = \frac{\sigma_{\max}}{\sigma_{\min}}
        = \frac{\sigma_1}{\sigma_r}
    \]
    It measures the \emph{ratio of maximum to minimum stretching}
    performed by $X$.
    A large $\kappa$ means the matrix nearly collapses some directions:
    small input perturbations in those directions become large output errors.
  \end{block}

  \begin{block}{Why $X^\top X$ Squares the Condition Number}
    Let $X = U\Sigma V^\top$ (SVD).
    Then $X^\top X = V\Sigma^2 V^\top$, whose eigenvalues are
    $\sigma_1^2 \geq \cdots \geq \sigma_r^2$.
    Therefore:
    \[
      \kappa(X^\top X)
        = \frac{\sigma_1^2}{\sigma_r^2}
        = \left(\frac{\sigma_1}{\sigma_r}\right)^{\!2}
        = \kappa(X)^2
    \]
    \textbf{Forming $X^\top X$ squares the condition number.}
    Example: if $\kappa(X) = 10^3$, then $\kappa(X^\top X) = 10^6$.
    If $\sigma_{\min}(X) = 10^{-8}$, then
    $\sigma_{\min}(X^\top X) = 10^{-16}$, which is at machine-epsilon:
    the computer may treat this eigenvalue as zero, destroying accuracy.
  \end{block}
\end{frame}

%------------------------------------------------------------------
\begin{frame}[shrink=30]{Summary of Basics --- Why PCA Uses SVD Directly}
  \begin{block}{Two Routes to the Principal Directions $V$}
    \textbf{Route 1 (Covariance / EVD):}
    \[
      X \;\xrightarrow{\text{form}}\; X^\top X
        \;\xrightarrow{\text{EVD}}\; V
    \]
    Eigenvalues of $X^\top X$ are $\lambda_i = \sigma_i^2$.
    Condition number goes from $\kappa(X)$ to $\kappa(X)^2$.

    \medskip

    \textbf{Route 2 (Direct SVD):}
    \[
      X \;\xrightarrow{\text{SVD}}\; U,\,\Sigma,\,V
    \]
    Modern algorithms (Golub--Kahan bidiagonalisation + QR iteration)
    apply orthogonal transformations (Householder reflections)
    \emph{directly} to $X$ --- they never form $X^\top X$.
    They compute $\sigma_i$ without squaring them, so numerical
    precision is preserved.
  \end{block}

  \begin{block}{Worked Condition Number Example}
    Suppose $\sigma_{\max}(X)=10^4,\ \sigma_{\min}(X)=10^{-4}$.

    \centering
    \renewcommand{\arraystretch}{1.4}
    \begin{tabular}{lcc}
      \toprule
      \textbf{Approach} & \textbf{Extremal values seen} & \textbf{$\kappa$}\\
      \midrule
      Direct SVD on $X$ & $10^4,\ 10^{-4}$ & $10^8$\\
      EVD on $X^\top X$ & $10^8,\ 10^{-8}$ & $10^{16}$\\
      \bottomrule
    \end{tabular}

    \medskip
    At $\kappa = 10^{16}$ the ratio exceeds double-precision
    machine epsilon ($\approx 10^{-16}$): the smallest eigenvalue
    is indistinguishable from zero numerically.
    That is why \texttt{sklearn.decomposition.PCA} uses SVD internally.
  \end{block}
\end{frame}

%------------------------------------------------------------------
% (C) AA^T EXAMPLE: EIGENVALUES AND EIGENVECTORS
%------------------------------------------------------------------
\begin{frame}[shrink=10]{Summary of Basics --- Computing Eigenvalues of $AA^\top$}
  \begin{block}{Setup}
    Let $A = \begin{bmatrix}3 & 0\\4 & 5\end{bmatrix}$.
    Compute $AA^\top$:
    \[
      AA^\top
        = \begin{bmatrix}3&0\\4&5\end{bmatrix}
          \begin{bmatrix}3&4\\0&5\end{bmatrix}
        = \begin{bmatrix}9&12\\12&41\end{bmatrix}
    \]
  \end{block}

  \begin{block}{Finding Eigenvalues via the Characteristic Polynomial}
    Solve $\det(AA^\top - \lambda I) = 0$:
    \[
      \det\begin{bmatrix}9-\lambda & 12\\12 & 41-\lambda\end{bmatrix}
        = (9-\lambda)(41-\lambda) - 144 = 0
    \]
    Expand: $\lambda^2 - 50\lambda + 369 - 144 = 0$,\;
    i.e.\ $\lambda^2 - 50\lambda + 225 = 0$.
    By the quadratic formula:
    \[
      \lambda = \frac{50 \pm \sqrt{2500 - 900}}{2}
              = \frac{50 \pm 40}{2}
      \quad\Rightarrow\quad
      \lambda_1 = 45,\quad \lambda_2 = 5
    \]
    Singular values of $A$: $\sigma_1 = \sqrt{45} = 3\sqrt{5}$,\;
    $\sigma_2 = \sqrt{5}$.
    Condition number: $\kappa(A) = \sqrt{45/5} = 3$.
  \end{block}
\end{frame}

%------------------------------------------------------------------
\begin{frame}[shrink=10]{Summary of Basics --- Eigenvectors of $AA^\top$ (Worked)}
  \begin{block}{Eigenvector for $\lambda_1 = 45$}
    $(AA^\top - 45I)\mathbf{u} = \mathbf{0}$:
    \[
      \begin{bmatrix}-36&12\\12&-4\end{bmatrix}
      \begin{bmatrix}x\\y\end{bmatrix}
      = \mathbf{0}
      \;\Rightarrow\; -36x + 12y = 0 \;\Rightarrow\; y = 3x
    \]
    Eigenvector (unnormalised): $[1,\ 3]^\top$.
    Normalised:
    $\mathbf{u}_1 = \frac{1}{\sqrt{10}}[1,\ 3]^\top$.
  \end{block}

  \begin{block}{Eigenvector for $\lambda_2 = 5$}
    $(AA^\top - 5I)\mathbf{u} = \mathbf{0}$:
    \[
      \begin{bmatrix}4&12\\12&36\end{bmatrix}
      \begin{bmatrix}x\\y\end{bmatrix}
      = \mathbf{0}
      \;\Rightarrow\; 4x + 12y = 0 \;\Rightarrow\; x = -3y
    \]
    Eigenvector (unnormalised): $[-3,\ 1]^\top$.
    Normalised:
    $\mathbf{u}_2 = \frac{1}{\sqrt{10}}[-3,\ 1]^\top$.
  \end{block}

  \begin{exampleblock}{Orthogonality Check}
    $\mathbf{u}_1^\top\mathbf{u}_2
     = \frac{1}{10}(1\cdot(-3) + 3\cdot 1) = 0$. \checkmark\;
    (As guaranteed by the Spectral Theorem: eigenvectors of a real
    symmetric matrix corresponding to \emph{distinct} eigenvalues
    are orthogonal.)
  \end{exampleblock}
\end{frame}

%------------------------------------------------------------------
% (D) BESSEL'S CORRECTION
%------------------------------------------------------------------
\begin{frame}[shrink=35]{Summary of Basics --- Bessel's Correction: The Key Lemma}
  \begin{block}{Claim}
    \[
      E\!\left[\sum_{i=1}^n (X_i - \bar{X})^2\right] = (n-1)\sigma^2
    \]
    where $X_1,\ldots,X_n \overset{\text{i.i.d.}}{\sim}$ distribution
    with mean $\mu$ and variance $\sigma^2$,
    and $\bar{X} = \frac{1}{n}\sum_{i=1}^n X_i$.
  \end{block}

  \begin{block}{Proof: Algebraic Decomposition}
    Write $X_i - \bar{X} = (X_i - \mu) - (\bar{X} - \mu)$.
    Then:
    \[
      \sum_{i=1}^n(X_i - \bar{X})^2
        = \sum_{i=1}^n(X_i-\mu)^2
          - 2(\bar{X}-\mu)\underbrace{\sum_{i=1}^n(X_i-\mu)}_{n(\bar{X}-\mu)}
          + n(\bar{X}-\mu)^2
    \]
    \[
      = \sum_{i=1}^n(X_i-\mu)^2 - n(\bar{X}-\mu)^2
    \]
    Take expectation:
    \[
      E\!\left[\sum_{i=1}^n(X_i-\bar{X})^2\right]
        = E\!\left[\sum_{i=1}^n(X_i-\mu)^2\right]
          - n\,E\!\left[(\bar{X}-\mu)^2\right]
    \]
    $= n\sigma^2 - n\cdot\dfrac{\sigma^2}{n} = n\sigma^2 - \sigma^2
    = (n-1)\sigma^2$.
    \hfill (using $E[(X_i-\mu)^2]=\sigma^2$ and
    $\mathrm{Var}(\bar{X})=\sigma^2/n$)
  \end{block}
\end{frame}

%------------------------------------------------------------------
\begin{frame}[shrink=45]{Summary of Basics --- Bessel's Correction: Bias and Unbiasedness}
  \begin{block}{Why Dividing by $n$ is Biased}
    Define $S_n^2 = \dfrac{1}{n}\sum_{i=1}^n(X_i - \bar{X})^2$.
    By the lemma and linearity of expectation ($E[cY] = cE[Y]$):
    \[
      E[S_n^2]
        = \frac{1}{n}\,E\!\left[\sum_{i=1}^n(X_i-\bar{X})^2\right]
        = \frac{1}{n}(n-1)\sigma^2
        = \frac{n-1}{n}\sigma^2
        \;\neq\; \sigma^2
    \]
    The estimator systematically \emph{underestimates} $\sigma^2$
    by a factor $\frac{n-1}{n}$.
  \end{block}

  \begin{block}{Bessel's Correction: Divide by $n-1$}
    Define $S^2 = \dfrac{1}{n-1}\sum_{i=1}^n(X_i-\bar{X})^2$.
    Then:
    \[
      E[S^2]
        = \frac{1}{n-1}\,(n-1)\sigma^2
        = \sigma^2
    \]
    $S^2$ is an \textbf{unbiased estimator} of $\sigma^2$.
  \end{block}

  \begin{block}{Degrees of Freedom: Why $n-1$?}
    The $n$ deviations $d_i = X_i - \bar{X}$ satisfy
    $\sum_{i=1}^n d_i = 0$ (since $\sum X_i = n\bar{X}$).
    This single linear constraint reduces the dimension of the
    space in which $(d_1,\ldots,d_n)$ lives from $n$ to $n-1$.
    Knowing any $n-1$ deviations determines the $n$-th:
    $d_n = -\sum_{i=1}^{n-1}d_i$.
    Hence there are $n-1$ \textbf{degrees of freedom} ---
    and we divide by $n-1$.
  \end{block}

  \begin{alertblock}{Connection to PCA}
    In PCA we use $\Sigma = \frac{1}{m}X_c^\top X_c$ (divide by $m$,
    not $m-1$) --- a population-variance convention common in ML.
    Eigenvectors are identical either way; only eigenvalue \emph{scale}
    differs by $\frac{m}{m-1}$, which is negligible for large $m$.
  \end{alertblock}
\end{frame}

\end{document}